$latexcv-preamble.tex()$

%-------------------------------------------------------------------------------
%	MINIMALISTIC
%
%	No filled bars and no rules between entries: headings are large, uppercase
%	and black, and everything else is separated by space alone. The accent
%	colour is spent on the entry titles rather than on furniture.
%-------------------------------------------------------------------------------

$if(geometry)$
\geometry{$for(geometry)$$geometry$$sep$,$endfor$}
$else$
\geometry{top=1.5cm, bottom=1.5cm, left=2cm, right=2cm}
$endif$

%-------------------------------------------------------------------------------
%	HEADINGS
%-------------------------------------------------------------------------------
\newcommand{\cvsection}[1]{%
  \vspace{6pt}
  \begin{flushleft}
    \LARGE\textcolor{black}{\uppercase{#1}}
  \end{flushleft}
  \vspace{-4pt}}

\newcommand{\cvsubsection}[1]{%
  \vspace{6pt}
  \begin{flushleft}
    \large\textcolor{bgcol}{#1}
  \end{flushleft}
  \vspace{-4pt}}

%-------------------------------------------------------------------------------
%	ENTRIES
%-------------------------------------------------------------------------------
% Stacked rather than tabular: the title on its own line in the accent colour,
% and the context beneath it in grey. Upstream sets the whole headline in a
% \vbox so that an entry is never broken from its first line.
\newcommand{\cventryhead}[4]{%
  \textcolor{sectcol}{\textbf{#2}}\\
  \textcolor{bgcol}{%
    \ifthenelse{\isempty{#3}}{}{#3}%
    \ifthenelse{\isempty{#4}}{}{\ifthenelse{\isempty{#3}}{}{, }\textit{#4}}%
    \ifthenelse{\isempty{#1}}{}{\ifthenelse{\isempty{#3}}{\ifthenelse{\isempty{#4}}{}{ | }}{ | }#1}}}

\newcommand{\vitaeentry}[5]{%
  \begin{flushleft}
    \vbox{\cventryhead{#1}{#2}{#3}{#4}\\[4pt]}
    #5
  \end{flushleft}
  \vspace{2pt}}

\newcommand{\vitaebrief}[4]{%
  \begin{flushleft}
    \vbox{\cventryhead{#1}{#2}{#3}{#4}}
  \end{flushleft}
  \vspace{3pt}}

\newcommand{\vitaeskill}[2]{%
  \noindent
  \begin{tabular*}{\linewidth}{@{}p{\cvskillwidth}@{}p{\dimexpr\linewidth-\cvskillwidth\relax}@{}}
    \textcolor{sectcol}{\textbf{#1}} & #2
  \end{tabular*}\par\vspace{4pt}}

% Plain dashes rather than chevrons or dots, in keeping with the rest.
\newenvironment{vitaewhy}{%
  \begin{itemize}[label=\textendash, topsep=0pt, partopsep=0pt, itemsep=2pt,
                  parsep=0pt, leftmargin=1.4em, labelsep=0.5em]%
}{%
  \end{itemize}%
}

% The disc the photograph is cut to, and the chord the rule below it sits on.
\newlength{\cvpicdiam}
\newlength{\cvpicrad}
\newlength{\cvpiccut}
\newlength{\cvpicfloor}

\begin{document}

%-------------------------------------------------------------------------------
%	PROFILE PICTURE (optional)
%-------------------------------------------------------------------------------
% Upstream's photograph is a disc whose foot the rule above the name cuts
% across, so that the head sits over the rule and the rest of the circle is
% never drawn. It gets that from its picture file -- a circle painted onto a
% white square, laid over the rule and then covered by the white box its name
% is set in. Neither half of that survives a photograph the document supplies:
% an ordinary square one would put its corners over the rule, and anything but
% a white page behind it would show the cover-up.
%
% So the shape is cut here instead of assumed. Two clips intersect: the disc,
% and the rectangle standing on the chord \cvpiccut above its foot. What is left
% is a circle with its bottom cut off, and the flat edge sits on the baseline
% for the rule to be drawn against.
$if(profilepic)$
\setlength{\cvpicdiam}{$if(profilepic-width)$$profilepic-width$$else$0.3\linewidth$endif$}
\setlength{\cvpicrad}{0.5\cvpicdiam}
% Upstream cuts 8% of the diameter off the foot of the disc.
\setlength{\cvpiccut}{0.08\cvpicdiam}
\setlength{\cvpicfloor}{\dimexpr\cvpiccut-\cvpicrad\relax}
{\parskip=0pt\relax\centering
 \begin{tikzpicture}[baseline=(current bounding box.south)]
   % Fixed first, so that the disc clipped inside it cannot grow the box back
   % down to its uncut foot.
   \useasboundingbox (-\cvpicrad,\cvpicfloor) rectangle (\cvpicrad,\cvpicrad);
   \clip (-\cvpicrad,\cvpicfloor) rectangle (\cvpicrad,\cvpicrad);
   \clip (0,0) circle (\cvpicrad);
   \node[inner sep=0pt] at (0,0)
     {\includegraphics[width=\cvpicdiam]{$profilepic$}};
 \end{tikzpicture}\par}
\nointerlineskip
$endif$

%-------------------------------------------------------------------------------
%	TITLE
%-------------------------------------------------------------------------------
% \color rather than \textcolor: \textcolor opens a paragraph for its argument,
% and the \parskip that comes with it would stand between the cut foot of the
% photograph and the rule that is meant to be drawn against it. With no
% photograph the rule is the first thing on the page, where that glue is
% discarded anyway, so nothing moves.
\begingroup\color{softcol}\hrule\endgroup
\vspace{10pt}
\begin{center}
  \HUGE \textsc{$name$$if(surname)$ $surname$$endif$}%
  $if(docname)$\cvtitlerule\textsc{$docname$}$endif$
\end{center}

$if(position)$
\begin{center}
  \small\textsc{$position$}
\end{center}
$endif$

%-------------------------------------------------------------------------------
%	CONTACT DETAILS
%-------------------------------------------------------------------------------
\vspace{4pt}
\begin{center}
  \footnotesize
  \renewcommand{\cvcontactsep}{\hspace{14pt}}
  \cvcontacts
\end{center}

\vspace{6pt}
\textcolor{softcol}{\hrule}
\vspace{4pt}

\normalsize

$latexcv-body.tex()$

$if(footer)$
\vspace{12pt}
\textcolor{softcol}{\hrule}
\vspace{4pt}
\begin{center}
  \footnotesize\textcolor{bgcol}{$footer$}
\end{center}
$endif$

\end{document}
